\documentclass[lang=cn,10pt]{elegantbook}

\title{清华附中初二上数学练习册}
\subtitle{知识点题库整理}

\author{齐齐爸}
\institute{珠江骏景课题组}
\date{April 9, 2026}
\version{4.3}
\bioinfo{自定义}{信息}

\extrainfo{不要以为抹消过去，重新来过，即可发生什么改变。—— 比企谷八幡}

\setcounter{tocdepth}{3}

\logo{logo-blue.png}
\cover{cover.jpg}

% 本文档命令
\usepackage{array}
\newcommand{\ccr}[1]{\makecell{{\color{#1}\rule{1cm}{1cm}}}}

% 添加 \usepackage{float}，以启用 [H] 选项。
\usepackage{float}

\usepackage{ctex}      % 中文支持
\usepackage{amsmath, amssymb}
\usepackage{graphicx}  % 图片插入

% ============================================================
% 章节题目公共排版命令
%
% 图片命令
% - \questionimage[<图片参数>]{<文件>}
%   用于“左文字右图片”结构里的右侧题图。
%   默认限制高度，避免题图过高。
% - \centerquestionimage[<图片参数>]{<盒子宽度>}{<文件>}
%   用于居中单图的底层命令。
% - \rightquestionimage[<图片参数>]{<盒子宽度>}{<文件>}
%   用于右对齐单图的底层命令。
% - \blankquestionimage[<宽度>]{<文件>}
% - \blankrightquestionimage[<宽度>]{<文件>}
% - \solutionquestionimage[<宽度>]{<文件>}
% - \solutionrightquestionimage[<宽度>]{<文件>}
%   按题型约定默认尺寸的公共命令，后续维护优先使用这一组。
%
% 图组命令
% - \begin{questionimagegroup} ... \end{questionimagegroup}
%   用于多图并排容器。
% - \questionimageitem[<图片参数>]{<盒子宽度>}{<文件>}{<图注>}
%   用于图组中的单张图片，图注可写 {} 表示留空。
% - \blankdoubleimageitem[<宽度>]{<文件>}{<图注>}
% - \blanktripleimageitem[<宽度>]{<文件>}{<图注>}
% - \solutiondoubleimageitem[<宽度>]{<文件>}{<图注>}
% - \solutiontripleimageitem[<宽度>]{<文件>}{<图注>}
%   用于填空题 / 解答题双图、三图的常用尺寸。
% - \questionimagegap[<间距>]
%   用于图组中图片间距，默认 0.02\textwidth。
%
% 题目列表命令
% - \questionlist[<起始题号>]  选择题外层列表
% - \blanklist[<起始题号>]     填空题外层列表
% - \solutionlist[<起始题号>]  解答题外层列表
% - \optionlist                选择题 A/B/C/D 选项列表
%
% 推荐尺寸
% - 选择题左文右图：右侧题图保持默认即可。
% - 填空题 / 解答题单图：建议优先使用 0.35\textwidth 左右。
% - 填空题 / 解答题双图：建议每张 0.32\textwidth 左右。
% - 填空题 / 解答题三图：建议每张 0.28\textwidth 左右。
%
% 模板库
%
% 选择题模板
% 1. 纯文字选择题：
%    \questionlist
%    \item 题干……
%    \optionlist
%    \item A
%    \item B
%    \item C
%    \item D
%    \end{enumerate}
%    \end{enumerate}
%
% 2. 左文字右图片选择题：
%    \questionlist
%    \item 题干……
%    \begin{questionwithimage}[width=0.9\linewidth]{foo.png}
%    \optionlist
%    \item A
%    \item B
%    \item C
%    \item D
%    \end{enumerate}
%    \end{questionwithimage}
%    \end{enumerate}
%
% 填空题模板
% 1. 从第 9 题开始的纯文字填空题：
%    \blanklist[9]
%    \item 第 9 题……
%    \item 第 10 题……
%    \item 第 11 题……
%    \end{enumerate}
%
% 2. 居中单图填空题：
%    \blanklist[9]
%    \item 题干……
%    \blankquestionimage{foo.png}
%    \end{enumerate}
%
% 3. 右对齐单图填空题：
%    \blanklist[9]
%    \item 题干……
%    \blankrightquestionimage{foo.png}
%    \end{enumerate}
%
% 4. 多图并排填空题：
%    \blanklist[9]
%    \item 题干……
%    \begin{questionimagegroup}
%    \blankdoubleimageitem{a.png}{图1}
%    \questionimagegap
%    \blankdoubleimageitem{b.png}{图2}
%    \end{questionimagegroup}
%    \end{enumerate}
%
% 5. 特殊高度限制的双图填空题：
%    \blanklist[9]
%    \item 题干……
%    \begin{questionimagegroup}
%    \questionimageitem[width=\linewidth,height=0.22\textheight,keepaspectratio]{0.28\textwidth}{a.png}{}
%    \questionimagegap[0.04\textwidth]
%    \questionimageitem[width=\linewidth,height=0.22\textheight,keepaspectratio]{0.28\textwidth}{b.png}{}
%    \end{questionimagegroup}
%    \end{enumerate}
%
% 6. 图注留空填空题：
%    \blanklist[9]
%    \item 题干……
%    \begin{questionimagegroup}
%    \blankdoubleimageitem{foo.png}{}
%    \end{questionimagegroup}
%    \end{enumerate}
%
% 解答题模板
% 1. 从第 15 题开始的纯文字解答题：
%    \solutionlist[15]
%    \item 证明题或计算题……
%    \end{enumerate}
%
% 2. 带居中单图的解答题：
%    \solutionlist[15]
%    \item 证明题或计算题……
%    \solutionquestionimage{foo.png}
%    \end{enumerate}
%
% 3. 带多图并排的解答题：
%    \solutionlist[15]
%    \item 证明题或计算题……
%    \begin{questionimagegroup}
%    \solutiondoubleimageitem{a.png}{图1}
%    \questionimagegap
%    \solutiondoubleimageitem{b.png}{图2}
%    \end{questionimagegroup}
%    \end{enumerate}
% ============================================================
\graphicspath{{figure/}{image/}{zhangz-book/figure/}{zhangz-book/image/}}

% 图片命令：默认参数
% `questionimage` 默认限制高度，`center/rightquestionimage` 默认按盒子宽度缩放。
\newcommand{\questionimagedefaultoptions}{width=\linewidth,height=0.22\textheight,keepaspectratio}
\newcommand{\centerquestionimagedefaultoptions}{width=\linewidth}
\newcommand{\blankquestionimagewidth}{0.35\textwidth}
\newcommand{\blankdoubleimagewidth}{0.32\textwidth}
\newcommand{\blanktripleimagewidth}{0.28\textwidth}
\newcommand{\solutionquestionimagewidth}{0.35\textwidth}
\newcommand{\solutiondoubleimagewidth}{0.32\textwidth}
\newcommand{\solutiontripleimagewidth}{0.28\textwidth}

% 图片命令：左文字右图片中的右侧题图
\newcommand{\questionimage}[2][]{%
  \vspace{0pt}%
  \begin{center}
    \if\relax\detokenize{#1}\relax
      \expandafter\includegraphics\expandafter[\questionimagedefaultoptions]{#2}%
    \else
      \includegraphics[#1]{#2}%
    \fi
  \end{center}
}

% 图片命令：左文字右图片题型
% 可选参数会透传给右侧图片，例如：
% \begin{questionwithimage}[width=0.85\linewidth,height=0.26\textheight,keepaspectratio]{foo.png}
%   ...
% \end{questionwithimage}
\newenvironment{questionwithimage}[2][]{%
  \def\questionwithimageoptions{#1}%
  \def\questionwithimagefile{#2}%
  \par\smallskip\noindent
  \begin{minipage}[t]{0.64\linewidth}
  \vspace{0pt}
}{%
  \end{minipage}
  \hfill
  \begin{minipage}[t]{0.32\linewidth}
  \vspace{0pt}%
  \begin{center}
    \ifx\questionwithimageoptions\empty
      \expandafter\includegraphics\expandafter[\questionimagedefaultoptions]{\questionwithimagefile}%
    \else
      \expandafter\includegraphics\expandafter[\questionwithimageoptions]{\questionwithimagefile}%
    \fi
  \end{center}
  \end{minipage}
}

% 图片命令：居中单图
% 参数顺序：图片选项、外层盒子宽度、图片文件。
\newcommand{\centerquestionimage}[3][]{%
  \begin{center}
  \begin{minipage}[b]{#2}
  \centering
  \if\relax\detokenize{#1}\relax
    \expandafter\includegraphics\expandafter[\centerquestionimagedefaultoptions]{#3}%
  \else
    \includegraphics[#1]{#3}%
  \fi
  \end{minipage}
  \end{center}
}

% 图片命令：右对齐单图
% 适合保留旧稿里“题干后右侧跟一张图”的阅读节奏。
\newcommand{\rightquestionimage}[3][]{%
  \begin{flushright}
  \begin{minipage}[b]{#2}
  \centering
  \if\relax\detokenize{#1}\relax
    \expandafter\includegraphics\expandafter[\centerquestionimagedefaultoptions]{#3}%
  \else
    \includegraphics[#1]{#3}%
  \fi
  \end{minipage}
  \end{flushright}
}

% 图片命令：按题型约定尺寸的单图包装
% 可选参数仅用于覆盖默认宽度；若需要高度限制等特殊参数，继续使用底层命令。
\NewDocumentCommand{\blankquestionimage}{O{\blankquestionimagewidth} m}{%
  \centerquestionimage[width=#1]{#1}{#2}%
}
\NewDocumentCommand{\blankrightquestionimage}{O{\blankquestionimagewidth} m}{%
  \rightquestionimage[width=#1]{#1}{#2}%
}
\NewDocumentCommand{\solutionquestionimage}{O{\solutionquestionimagewidth} m}{%
  \centerquestionimage[width=#1]{#1}{#2}%
}
\NewDocumentCommand{\solutionrightquestionimage}{O{\solutionquestionimagewidth} m}{%
  \rightquestionimage[width=#1]{#1}{#2}%
}

% 图组命令：多图并排容器
\newenvironment{questionimagegroup}{%
  \begin{center}
}{%
  \end{center}
}

% 图组命令：图组中的单张图片
% 图注可传空字符串 {}，表示只显示图片不显示图注。
\newcommand{\questionimageitem}[4][]{%
  \begin{minipage}[b]{#2}
  \centering
  \if\relax\detokenize{#1}\relax
    \expandafter\includegraphics\expandafter[\centerquestionimagedefaultoptions]{#3}%
  \else
    \includegraphics[#1]{#3}%
  \fi
  \if\relax\detokenize{#4}\relax\else
    \par\smallskip
    {\centering\small #4\par}
  \fi
  \end{minipage}
}

% 图组命令：按题型约定尺寸的双图、三图包装
\NewDocumentCommand{\blankdoubleimageitem}{O{\blankdoubleimagewidth} m m}{%
  \questionimageitem{#1}{#2}{#3}%
}
\NewDocumentCommand{\blanktripleimageitem}{O{\blanktripleimagewidth} m m}{%
  \questionimageitem{#1}{#2}{#3}%
}
\NewDocumentCommand{\solutiondoubleimageitem}{O{\solutiondoubleimagewidth} m m}{%
  \questionimageitem{#1}{#2}{#3}%
}
\NewDocumentCommand{\solutiontripleimageitem}{O{\solutiontripleimagewidth} m m}{%
  \questionimageitem{#1}{#2}{#3}%
}

% 图组命令：图组间距
\newcommand{\questionimagegap}[1][0.02\textwidth]{%
  \hspace{#1}%
}

% 题目列表命令
% 可选参数表示起始题号，例如 \blanklist[9]。
\newcommand{\questionlist}[1][1]{%
  \begin{enumerate}[label=\arabic*., leftmargin=*, start=#1, itemsep=0.5\baselineskip]
}
\newcommand{\blanklist}[1][1]{%
  \begin{enumerate}[label=\arabic*., leftmargin=*, start=#1, itemsep=0.5\baselineskip]
}
\newcommand{\solutionlist}[1][1]{%
  \begin{enumerate}[label=\arabic*., leftmargin=*, start=#1, itemsep=0.8\baselineskip]
}
\newcommand{\optionlist}{%
  \begin{enumerate}[label=(\Alph*), itemsep=0.2\baselineskip]
}
\makeatletter
% 兼容从项目根目录或 zhangz-book 目录发起编译。
\def\input@path{{./}{zhangz-book/}}
\makeatother


% \usepackage{wrapfig}

% 修改标题页的橙色带
% \definecolor{customcolor}{RGB}{32,178,170}
% \colorlet{coverlinecolor}{customcolor}

\begin{document}

\maketitle
\frontmatter

\tableofcontents

\mainmatter


%第十二章 平面直角坐标系与几何变换

% 12.1 平面直角坐标系

% 12.2 用坐标表示位置

% 12.3 平面直角坐标系中的基本公式

% 12.4 平移

\chapter{12.4 平移}

\section*{一、选择题（共 8 小题）}

\questionlist
\item
  \begin{questionwithimage}{12.4-1.png}
  如图，$\triangle ABC$ 经过怎样的平移得到 $\triangle DEF$（\quad）
  \optionlist
  \item 把$\triangle ABC$ 向左平移 4 个单位，再向下平移 2 个单位
  \item 把$\triangle ABC$ 向右平移 4 个单位，再向下平移 2 个单位
  \item 把$\triangle ABC$ 向右平移 4 个单位，再向上平移 2 个单位
  \item 把$\triangle ABC$ 向左平移 4 个单位，再向上平移 2 个单位
  \end{enumerate}
  \end{questionwithimage}

\item 在平面内，将一个图形沿某个方向移动一定距离，这样的图形变换为平移，如图，将网格中的三条线段沿网格线的方向（水平或垂直）平移后组成一个首尾依次相接的三角形，至少需要移动（\quad）
  \begin{questionwithimage}{12.4-2.png}
  \optionlist
  \item 12 格
  \item 11 格
  \item 9 格
  \item 8 格
  \end{enumerate}
  \end{questionwithimage}

\item 如图，图案⑥是由①②③④⑤五种基本图形中的两种拼接而成的，这两种基本图形是（\quad）
  \begin{questionwithimage}{12.4-3.png}
  \optionlist
  \item ①⑤
  \item ②③
  \item ④⑥
  \item ②④
  \end{enumerate}
  \end{questionwithimage}

\item 如图所示，半圆 $AB$ 平移到半圆 $CD$ 的位置时所扫过的面积为（\quad）
  \begin{questionwithimage}{12.4-4.png}
  \optionlist
  \item 3
  \item $3+\pi$
  \item 6
  \item $6+\pi$
  \end{enumerate}
  \end{questionwithimage}

\item 如图，在长方形 $ABCD$ 中，$AB=8$，$BC=5$，则图中四个小长方形的周长和为（\quad）
  \begin{questionwithimage}{12.4-5.png}
  \optionlist
  \item 13
  \item 23
  \item 24
  \item 26
  \end{enumerate}
  \end{questionwithimage}

\item 如图，相邻两线段互相垂直，两只蜗牛均同时从点出发到达点，蜗牛甲沿着“$A\rightarrow B\rightarrow C$”路线走，蜗牛乙沿着“$A\rightarrow D\rightarrow E\rightarrow F\rightarrow G\rightarrow H\rightarrow I\rightarrow J\rightarrow C$”的路线走，若他们的爬行速度相同，则先到达点 C 的是（\quad）
  \par\smallskip\noindent
  \begin{minipage}[t]{0.54\linewidth}
  \vspace{0pt}
  \optionlist
  \item 蜗牛甲
  \item 蜗牛乙
  \item 同时到达
  \item 无法确定
  \end{enumerate}
  \end{minipage}
  \hfill
  \begin{minipage}[t]{0.32\linewidth}
  \questionimage{12.4-6.png}
  \end{minipage}
\end{enumerate}

% 以下题目根据截图转录，图形暂未单独拆分为本地图片文件。
\questionlist[7]
\item 如图，将$\triangle ABC$沿$BC$方向平移3 cm得到$\triangle DEF$，若$\triangle ABC$的周长为16 cm，则四边形$ABFD$的周长为（\quad）【图不清楚，可忽略】
  \optionlist
  \item 16 cm
  \item 22 cm
  \item 20 cm
  \item 24 cm
  \end{enumerate}

\item 如图，由$\triangle ABC$平移可得到的三角形有几个（\quad）【图不清楚，可忽略】
  \optionlist
  \item 3个
  \item 5个
  \item 6个
  \item 7个
  \end{enumerate}
\end{enumerate}

\section*{二、填空题（共 6 小题）}

\blanklist[9]

\item 在如图所示的长方体中，可以由线段$AB$平移而得到的线段有\underline{\hspace{3cm}}条。

\centerquestionimage[width=0.55\textwidth]{0.35\textwidth}{12.4-9.png}

\item 如图，在长为$a$米、宽为$b$米的长方形耕地上修筑宽度均为1米的道路（如阴影部分），其余空白部分用于种植草坪，则草坪面积为\underline{\hspace{4cm}}米$^2$。

\centerquestionimage[width=0.55\textwidth]{0.35\textwidth}{12.4-10.png}

\item 在等腰$\triangle ABC$中，$AB=AC$，则$BC$边上的中线、高线和$\angle BAC$的平分线重合于$AD$（如图一）。若将等腰$\triangle ABC$的顶点$A$向右平行移动后，得到$\triangle A'BC$（如图二），那么，此时$BC$边上的中线、$BC$边上的高线和$\angle BA'C$的平分线应依次分别是\underline{\hspace{1.4cm}}、\underline{\hspace{1.4cm}}、\underline{\hspace{1.4cm}}。（填$A'D$、$A'E$、$A'F$）

\centerquestionimage[width=0.6\textwidth]{0.5\textwidth}{12.4-11.png}

\item 如图，把$\angle AOB$沿着直线$MN$平移一定的距离，得到$\angle CPD$。若$\angle AOM=40^\circ$，$\angle DPN=40^\circ$，则$\angle AOB=$\underline{\hspace{2cm}}。

\centerquestionimage[width=0.35\textwidth]{0.55\textwidth}{12.4-12.png}

\item 边长为2的等边$\triangle ABC$与等边$\triangle DEF$互相重合，将$\triangle ABC$沿直线$l$向右平移$m$个单位长度，将$\triangle DEF$向右也平移$n$个单位长度，如图，当$C$、$E$是线段$BF$的三等分点时，$m$的值为\underline{\hspace{2cm}}。

\centerquestionimage[width=0.85\textwidth]{0.35\textwidth}{12.4-13.png}


\item 如图，长方形 $ABCD$ 中，$AB=6$，第一次平移长方形 $ABCD$ 沿 $AB$ 的方向向右平移 5 个单位，得到长方形 $A_1B_1C_1D_1$，第二次平移将长方形 $A_1B_1C_1D_1$ 沿 $A_1B_1$ 的方向向右平移 5 个单位，得到长方形 $A_2B_2C_2D_2$，...，第 $n$ 次平移将长方形 $A_{n-1}B_{n-1}C_{n-1}D_{n-1}$ 沿 $A_{n-1}B_{n-1}$ 的方向平移 5 个单位，得到长方形 $A_nB_nC_nD_n$ ($n \geq 2$)，若 $A_nB_n$ 的长度为 56，则 $n=\underline{\quad }$。

% 插入图片：此处应有第14题图


\vspace{0.2cm}
\centerquestionimage[width=0.90\textwidth]{0.35\textwidth}{12.4-14.png}


\end{enumerate}

\section*{三、解答题（共 5 小题）}


\solutionlist[15]


\item 如图，将 $Rt\triangle ABC$ 沿 $AB$ 方向平移得到 $Rt\triangle DEF$，已知 $BE=6$，$EF=8$，$CG=3$，求阴影部分的面积。

% 第15题图
\rightquestionimage[width=0.55\textwidth]{0.55\textwidth}{12.4-15.png}

\vspace{1cm}

\item 如图，$\triangle A_1B_1C_1$ 是 $\triangle ABC$ 向右平移 4 个单位长度后得到的，且三个顶点的坐标分别为 $A_1(1, 1)$，$B_1(4, 2)$，$C_1(3, 4)$。

(1) 请画出 $\triangle ABC$，并写出点 $A$，$B$，$C$ 的坐标；

(2) 求出 $\triangle AOA_1$ 的面积。

% 第16题图
\rightquestionimage[width=0.75\textwidth]{0.35\textwidth}{12.4-16.png}

\vspace{2cm}

\item 如图，在直角坐标系 $xOy$ 中，$A(-1, 0)$，$B(3, 0)$，将 $A, B$ 同时分别向上平移 2 个单位，再向右平移 1 个单位，得到的对应点分别为 $D, C$，连接 $AD, BC$。

(1) 直接写出点 $C, D$ 的坐标：$C$ \_\_\_\_\_\_，$D$ \_\_\_\_\_\_；

(2) 四边形 $ABCD$ 的面积为 \_\_\_\_\_\_；

(3) $P$ 在线段 $BC$ 上，$\triangle OPD$ 面积是 $\frac{11}{3}$，求 $P$ 点的坐标。

% 第17题图
\rightquestionimage[width=0.75\textwidth]{0.35\textwidth}{12.4-17.png}

\vspace{3cm}


\item 如图，长方形 $ABCD$ 在坐标平面内，点 $A$ 的坐标是 $A(2, 1)$，且边 $AB, CD$ 与 $x$ 轴平行，边 $AD, BC$ 与 $y$ 轴平行，点 $B, C$ 的坐标分别为 $B(a, 1)$，$C(a, c)$，且 $a, c$ 满足关系式：$c = \sqrt{a-6} + \sqrt{6-a} + 3$。


(1) 求 $B, C, D$ 三点的坐标；

(2) 怎样平移，才能使 $A$ 点与原点重合？平移后点 $B, C, D$ 的对应分别为 $B_1C_1D_1$，求四边形 $OB_1C_1D_1$ 的面积；

(3) 平移后在 $x$ 轴上是否存在点 $P$，连接 $PD$，使 $S_{\triangle COP} = S_{\text{四边形} OBCD}$？若存在这样的点 $P$，求出点 $P$ 的坐标；若不存在，试说明理由。

% 第18题图
\rightquestionimage[width=0.85\textwidth]{0.5\textwidth}{12.4-18.png}

\vspace{3cm}

\item 我们约定，若一个三角形（记为 $\triangle A_1$）是由另一个三角形（记为 $\triangle A$）通过一次平移，或绕其任一边的中点旋转 $180^\circ$ 得到的，则称 $\triangle A_1$ 是由 $\triangle A$ 复制的。以下的操作中每一个三角形可以复制一次，复制过程可以一直进行下去。如图 1 是由 $\triangle A$ 复制出 $\triangle A_1$，又由 $\triangle A_1$ 复制出 $\triangle A_2$，再由 $\triangle A_2$ 复制出 $\triangle A_3$，形成了一个大三角形，记作 $\triangle B$。以下各题中的复制均是由 $\triangle A$ 开始的，由复制形成的多边形中的任意两个小三角形（指与 $\triangle A$ 全等的三角形）之间既无缝隙也无重叠。


(1) 图 1 中标出的是一种可能的复制结果，它用到\_\_\_\_\_\_次平移，\_\_\_\_\_\_次旋转。若由复制形成的 $\triangle C$ 的一条边上有 11 个小三角形（指有一条边在该边上的小三角形），则 $\triangle C$ 中含有\_\_\_\_\_\_个小三角形；

(2) 若 $\triangle A$ 是正三角形，你认为通过复制能形成的正多边形是\_\_\_\_\_\_。

(3) 在复制形成四边形的过程中，小明用到了两次平移一次旋转，你能用两次旋转一次平移复制形成一个四边形吗？如果能，请在图 2 的方框内画出草图，并仿照图 1 作出标记；如果不能，请说明理由；

(4) 图 3 是正五边形 $EFGHI$，其中心是 $O$。连接 $O$ 点与各顶点。将其中的一个三角形记为 $\triangle A$，小明认为正五边形 $EFGHI$ 是由复制形成的一种结果，你认为他的说法对吗？请判断并说明理由。

% 第19题图
\rightquestionimage[width=0.7\textwidth]{0.7\textwidth}{12.4-19.png}


\end{enumerate}



\chapter{12.5 用坐标表示平移}


\subsection*{一、选择题（共 8 小题）}

\questionlist

\item 在平面直角坐标系中，若将三角形上各点的纵坐标都减去 3，横坐标保持不变，则所得图形在原图形基础上（\quad ）

  \optionlist
  \item 向左平移了 3 个单位
  \item 向下平移了 3 个单位
  \item 向上平移了 3 个单位
  \item 向右平移了 3 个单位
  \end{enumerate}

\item 在平面直角坐标系中，$\triangle ABC$ 的三个顶点坐标分别为 $A(4, 5)$，$B(1, 2)$，$C(4, 2)$，将 $\triangle ABC$ 向左平移 5 个单位后，$A$ 的对应点 $A_1$ 的坐标是（\quad）

  \optionlist
  \item (0, 5)
  \item (-1, 5)
  \item (9, 5)
  \item (-1, 0)
  \end{enumerate}

\item 点 $A(4, 3)$ 经过平移后得点 $B(6, -3)$，它的平移过程是（\quad）

  \optionlist
  \item 向右平移 2 个单位后再向下平移 6 个单位
  \item 向左平移 2 个单位后再向下平移 2 个单位
  \item 向左平移 2 个单位后再向上平移 6 个单位
  \item 向右平移 6 个单位后再向上平移 2 个单位
  \end{enumerate}

\item 在平面直角坐标系中，已知点 $O(0, 0)$，$A(2, 4)$，将线段 $OA$ 沿 $x$ 轴向左平移 2 个单位，再向上平移 2 个单位，得到 $O, A$ 两点的对应点分别为 $O_1, A_1$，则 $O_1, A_1$ 的坐标分别是（\quad）

  \optionlist
  \item (-2, 2), (4, 6)
  \item (-2, 2), (0, 6)
  \item (2, -2), (4, 6)
  \item (2, -2), (0, 6)
  \end{enumerate}

\item 在平面直角坐标系中，将点 $A(x, y)$ 向左平移 5 个单位长度，再向上平移 3 个单位长度后与点 $B(-3, 2)$ 重合，则点 $A$ 的坐标是（\quad）

  \optionlist
  \item (2, 5)
  \item (-8, 5)
  \item (-8, -1)
  \item (2, -1)
  \end{enumerate}

\item 在 $\triangle ABC$ 内任意一点 $P(a, b)$ 经过平移后对应点 $P_1(c, d)$，已知 $A(3, 2)$ 在经过此次平移后对应点 $A_1$ 的坐标为 $(5, -1)$，则 $a+b-c-d$ 的值为（\quad）

  \optionlist
  \item -5
  \item -1
  \item 1
  \item 5
  \end{enumerate}

\item 在坐标平面上两点 $A(-a+2, -b+1)$，$B(3a, b)$，若点 $A$ 向右移动 2 个单位长度后，再向下移动 3 个单位长度后与点 $B$ 重合，则点 $B$ 所在的象限为（\quad）

  \optionlist
  \item 第一象限
  \item 第二象限
  \item 第三象限
  \item 第四象限
  \end{enumerate}






\item 图，已知直角坐标系中的点 $A, B$ 的坐标分别为 $A(2, 4)$，$B(4, 0)$，且 $P$ 为 $AB$ 的中点。若将线段 $AB$ 向右平移 3 个单位后，与点 $P$ 对应的点为 $Q$，则点 $Q$ 的坐标是（\quad）
  \par\smallskip\noindent
  \begin{minipage}[t]{0.44\linewidth}
  \vspace{0pt}
 \optionlist
  \item (3, 2)
  \item (6, 2)
  \item (6, 4)
  \item (3, 5)
  \end{enumerate}
  \end{minipage}
  \hfill
  \begin{minipage}[t]{0.2\linewidth}
  \questionimage{12.5-8.png}
  \end{minipage}
  


\end{enumerate}

\subsection*{二、填空题（共 6 小题）}

\blanklist[9]

\item 如图，将线段 $AB$ 平移，使 $B$ 点到 $C$ 点，则平移后 $A$ 点的坐标为\_\_\_\_\_\_。

% 插入图片：此处应有第9题图
\rightquestionimage[width=0.35\textwidth]{0.35\textwidth}{12.5-9.png}


\item 在平面直角坐标系中，有一条线段 $AB$，已知点 $A(-2, 0)$ 和 $B(0, 2)$，平移线段 $AB$ 得到线段 $A_1B_1$。若点 $A$ 的对应点 $A_1$ 的坐标为 $(1, 3)$，则线段 $A_1B_1$ 的中点坐标是\_\_\_\_\_\_。
\vspace{1cm}

\item 已知点 $P(-4, 6-3b)$，先向左平移 2 个单位，再向下平移 3 个单位，恰好落在 $x$ 轴的负半轴上，则 $b$ 应为\_\_\_\_\_\_。
\vspace{2cm}
\item 在平面直角坐标系中，规定一个点先向上平移 2 个单位，再向右平移 1 个单位为 1 次运动。点 $(-2, -3)$ 经过\_\_\_\_\_\_次这样的运动后到达点 $P'(4, 9)$。
\vspace{1cm}
\item 在平面直角坐标系中，孔明玩走棋的游戏，其走法是：棋子从原点出发，第 1 步向右走 1 个单位长度，第 2 步向右走 2 个单位长度，第 3 步向上走 1 个单位长度，第 4 步向右走 1 个单位长度……依此类推，第 $n$ 步的走法是：当 $n$ 能被 3 整除时，则向上走 1 个单位长度；当 $n$ 能被 3 整除时，则向下走 1 个单位长度；当 $n$ 能被 3 整除时，若 $n$ 能被 3 整除，则向右走 2 个单位长度，若 $n$ 能被 3 整除，则向左走 2 个单位长度。当走完第 100 步时，棋子所在位置的坐标是\_\_\_\_\_\_。
\vspace{2cm}
\item 初一年级某班有 54 名学生，所在教室有 6 行 9 列座位，用 $(m, n)$ 表示第 $m$ 行第 $n$ 列的位置，新学期准备调整座位，设某个学生原来的座位为 $(m, n)$，如果调整后的座位为 $(i, j)$，则称该生作了平移 $(a, b) = [m-i, n-j]$，并称 $a+b$ 为该生的位置数。若某生的位置数为 10，则当 $m+n$ 取最小值时，$m+n$ 的最大值为\_\_\_\_\_\_。
\vspace{2cm}

\end{enumerate}

\subsection*{三、解答题（共 4 小题）}

\solutionlist[15]

\item 五边形 $ABCDE$ 的顶点坐标分别为 $A(0, 6)$、$B(-3, -3)$、$C(-1, 0)$、$D(1, 0)$、$E(3, 3)$，将五边形 $ABCDE$ 看成经过一次平移后得到 $A_1B_1C_1D_1E_1$，其中顶点 $A$ 的对应点是 $A_1(-3, 10)$。

(1) 请写出其它对应点的坐标；

(2) 请写出这一平移的平移方向和平移距离。

\vspace{4cm}

\item 已知点 $P(2a-12, 1-a)$ 位于第三象限，点 $Q(x, y)$ 位于第二象限且是由点 $P$ 向上平移一定单位长度得到的。

(1) 若点 $P$ 的纵坐标为 $-3$，试求出 $a$ 的值；

(2) 在 (1) 的条件下，试求出符合条件的点 $Q$ 的坐标。

若点 $P$ 的横、纵坐标都是整数，试求出 $a$ 的值以及线段 $PQ$ 长度的取值范围。

\rightquestionimage[width=0.35\textwidth]{0.35\textwidth}{12.5-16.png}

\vspace{5cm}


\item 图中的图形是将坐标 $(0, 0), (1, 0), (3, 0), (2, 1), (3, 4), (5, 3), (5, 2), (3, 2)$ 的点用线段依次连接而成的。

% 插入图片：此处应有第17题图

下面将以上各点做如下变化：

(1) 横坐标保持不变，纵坐标分别乘以 $-1$，所得图案与原图案有什么变化？

(2) 横坐标和纵坐标都乘以 $-1$，所得图案与原图案相比有什么变化？

(3) 横坐标加 $1$，纵坐标加 $2$，所得图案与原图案相比有什么变化？

\rightquestionimage[width=0.35\textwidth]{0.35\textwidth}{12.5-17.png}

\vspace{5cm}

\newpage 

\item 操作与探究：

(1) 对数轴上的点 $P$ 进行如下操作：先把点 $P$ 表示的数乘以 $\frac{1}{3}$，再把所得数对应的点向右平移 $1$ 个单位，得到点 $P'$ 的对应点 $P'$。

点 $A, B$ 在数轴上，对线段 $AB$ 上的每个点进行上述操作后得到线段 $A'B'$，其中点 $A, B$ 的对应点分别为 $A', B'$。如图1，若点 $A$ 表示的数是 $-3$，则点 $A'$ 表示的数是\_\_\_\_\_\_；若点 $B'$ 表示的数是 $2$，则点 $B$ 表示的数是\_\_\_\_\_\_；已知线段 $AB$ 上的点 $E$ 经过上述操作后得到的对应点 $E'$ 与点 $E$ 重合，则点 $E$ 表示的数是\_\_\_\_\_\_。

% 插入图片：此处应有图1

(2) 如图2，在平面直角坐标系 $xOy$ 中，对正方形 $ABCD$ 及其内部的每个点进行如下操作：把每个点的横、纵坐标都乘以同一个实数 $a$，将得到的点先向右平移 $m$ 个单位，再向上平移 $n$ 个单位（$m > 0, n > 0$），得到正方形 $A'B'C'D'$ 及其内部的点，其中点 $A, B$ 的对应点分别为 $A', B'$。已知正方形 $ABCD$ 内部的一个点 $F$ 经过上述操作后得到的对应点 $F'$ 与点 $F$ 重合，求点 $F$ 的坐标。

% 插入图片：此处应有图2
\rightquestionimage[width=0.35\textwidth]{0.35\textwidth}{12.5-18.png}

\vspace{5cm}

\end{enumerate}


\chapter{12.6 轴对称与轴对称图形}


\section*{二、填空题（共 6 小题）}

\blanklist[9]

\item 线段是轴对称图形，它有\_\_\_\_\_\_条对称轴。

\item 将一根细绳对折 4 次后，用剪刀从中间剪断，这根细绳共被剪成了\_\_\_\_\_\_段。

\item 小明从镜子里看到镜子对面电子钟的像如图所示：
\[
00:25:31
\]
实际时间是\_\_\_\_\_\_。

% 插入图片：此处应有镜子中的电子钟像图

\item 如图，光线 $a$ 照射到平面镜 $CD$ 上，然后在平面镜 $AB$ 和 $CD$ 之间来回反射，这时光线的入射角等于反射角。若已知 $\angle 1 = 50^\circ$，$\angle 2 = 55^\circ$，则 $\angle 3 =$\_\_\_\_\_\_。

% 插入图片：此处应有第12题图

% \item 如图，$OE$ 是 $\angle AOB$ 的平分线，$BD \perp OA$ 于点 $D$，$AC \perp BO$ 于点 $C$，则关于直线 $OE$ 对称的三角形共有\_\_\_\_\_\_对。

% 插入图片：此处应有第13题图

% \item 如图，在平面直角坐标系中有三点：$A(1, 0)$，$B(-2, 0)$，$C(0, 2)$，点 $D$ 为 $y$ 轴上一点，当四边形 $ACBD$ 为轴对称图形时，点 $D$ 的坐标是\_\_\_\_\_\_。

% 插入图片：此处应有第14题图

\end{enumerate}

\rightquestionimage[width=0.4\textwidth]{0.4\textwidth}{12.6-12.png}

\section*{三、解答题（共 6 小题）}

\solutionlist[15]

  \item $\triangle ABC$ 的三边长分别为：$AB = 2a^2 - a - 7$，$BC = 10 - a^2$，$AC = a$。

(1) 求 $\triangle ABC$ 的周长（请用含有 $a$ 的代数式表示）；

(2) 当 $a = 2.5$ 和 $3$ 时，三角形都存在吗？若存在，求出 $\triangle ABC$ 的周长；若不存在，请说明理由；

(3) 若 $\triangle ABC$ 与 $\triangle DEF$ 成轴对称图形，其中点 $A$ 与点 $D$ 是对称点，点 $B$ 与点 $E$ 是对称点，$EF = 2$，$DF = 3 - b$，求 $a - b$ 的值。

\vspace{4cm}

\item 设 $l_1$ 和 $l_2$ 是两平行相对的镜子，如果把一个小球放在 $l_1$ 和 $l_2$ 之间（如图），试问：

\rightquestionimage[width=0.4\textwidth]{0.4\textwidth}{12.6-17.png}

\vspace{5cm}

(1) 小球 $A$ 在镜 $l_1$ 中的像 $A'$ 在什么位置？

(2) 小球 $A$ 在镜 $l_2$ 中的像 $A'$ 在镜 $l_2$ 中的像 $A''$ 又在什么位置？分别画在图上；

(3) 小球 $A$ 和像 $A''$ 之间的距离与 $l_1$ 和 $l_2$ 之间的距离有什么关系？

\item 如图，选择适当的方向击打白球，可使白球经过两次反弹后将红球撞入底袋。白球在运动过程中，$\angle 1 = \angle 2$，$\angle 3 = \angle 4$。如果红球与洞口的连线与台球桌边缘的夹角 $\angle 5 = 25^\circ$，那么选择 $\angle 1$ 是多少度，才能保证红球能直接入袋？为什么？


\rightquestionimage[width=0.4\textwidth]{0.4\textwidth}{12.6-18.png}


\vspace{5cm}

\item 下列图中是由字母 A 和 H 构成的（把 A、H 视为轴对称图形）。

% 插入图片：此处应有第1题图（由字母 A 和 H 构成的轴对称图形图案）
\centerquestionimage[width=0.35\textwidth]{0.35\textwidth}{12.6-19.png}


(1) 仔细观察其中的变化规律，回答下列问题：

① 第 100 个字母是什么？

② 图形中的字母 A 在前 2025 个字母中一共出现多少次？

(2) 从左往右在图案中至少取多少个（多于 1 个）字母能构成一次轴对称？字母个数为多少个（多于 1 个）字母能构成轴对称？

\vspace{4cm}

\item 如图，平面直角坐标系中，$A(0, x)$、$B(y, 0)$、$C(z, 0)$，在 $B$、$C$ 两点各有一个平面镜，其中在 $B$ 点的平面镜沿 $x$ 轴方向，从 $P$ 点发射两条光线 $PA$、$PB$，反射光线 $BD$ 经 $A$ 点和反射光线 $CD$ 相交。

\rightquestionimage[width=0.4\textwidth]{0.4\textwidth}{12.6-20.png}

(1) 若 $x$、$y$、$z$ 满足 $(2x+y-1)^2 + |y+z-1| = -(z-2)^2$，求 $\triangle ABC$ 的面积；

(2) 若两条入射光线 $PA$、$PB$ 的夹角（$\angle BPC$）为 $28^\circ$，要想让两条反射光线 $BD$、$CD$ 的夹角（$\angle BDC$）为 $36^\circ$，问平面镜 $MN$ 与 $x$ 轴夹角的度数。

\vspace{4cm}

\end{enumerate}


\chapter{12.7 线段的垂直平分线-1}


\section*{三、解答题（共 4 小题）}

\solutionlist[15]

\item 作图题：（不写作法，但必须保留作图痕迹）

如图：某地有两所大学和两条相交的公路，（点 $M, N$ 表示大学，$AO, BO$ 表示公路）。现计划修建一座物资仓库，希望仓库到两所大学的距离相等，到两条公路的距离也相等。你能确定仓库 $P$ 应该建在什么位置吗？在所给的图形中画出你的设计方案。

\rightquestionimage[width=0.4\textwidth]{0.4\textwidth}{12.7-15.png}

\vspace{3cm}

\item 在 $\triangle ABC$ 中，$AD$ 是高，在线段 $DC$ 上取一点 $E$，使 $BD = DE$，已知 $AB + BD = DC$，求证：$E$ 点在线段 $AC$ 的垂直平分线上。

\rightquestionimage[width=0.4\textwidth]{0.4\textwidth}{12.7-16.png}

\vspace{3cm}

\newpage

\item $\triangle ABC$ 中，$DE, FG$ 分别垂直平分边 $AB, AC$，垂足分别为点 $D, G$。

\rightquestionimage[width=0.4\textwidth]{0.4\textwidth}{12.7-17.png}

(1) 如图，①若 $\angle B = 30^\circ$，$\angle C = 40^\circ$，求 $\angle EAF$ 的度数；

②如果 $BC = 10$，求 $\triangle EAF$ 的周长；

③若 $AE \perp AF$，则 $\angle BAC = \underline{\quad }^\circ$。

(2) 若 $\angle BAC = n^\circ$，则 $\angle EAF = \underline{\quad }^\circ$。（用含 $n$ 代数式表示）

\vspace{7cm}


\item 如图，在 $\triangle ABC$ 中，$AB$ 边的垂直平分线 $l_1$ 交 $BC$ 于点 $D$，$AC$ 边的垂直平分线 $l_2$ 交 $BC$ 于点 $E$，$l_1$ 与 $l_2$ 相交于点 $O$，连接 $OB$、$OC$，若 $\triangle ADE$ 的周长为 $6\text{cm}$，$\triangle OBC$ 的周长为 $16\text{cm}$。

% 插入图片：此处应有第18题图

1. 求线段 $BC$ 的长；

2. 连接 $OA$，求线段 $OA$ 的长；

3. 若 $\angle BAC = 120^\circ$，求 $\angle DAE$ 的度数。

\rightquestionimage[width=0.4\textwidth]{0.4\textwidth}{12.7-18.png}


\end{enumerate}

\chapter{12.7 线段的垂直平分线-2}


\section*{三、解答题（共 4 小题）}

\solutionlist[15]

\item 如图，$\angle ACB = 90^\circ$，$AC = BC$，$D$ 为 $\triangle ABC$ 外一点，且 $AD = BD$，$DE \perp AC$ 交 $CA$ 的延长线于 $E$ 点。求证：$DE = AE + BC$。


\rightquestionimage[width=0.35\textwidth]{0.35\textwidth}{12.7-2-17.png}

\vspace{3cm}




\end{enumerate}


\chapter{12.8 坐标系中的轴对称}


\section*{一、选择题（共8小题）}


\questionlist

\item 如果点 $P(-2, b)$ 和点 $Q(a, -3)$ 关于 $x$ 轴对称，则 $a+b$ 的值是（\quad）

  \optionlist
  \item -1
  \item 1
  \item -5
  \item 5
  \end{enumerate}

\item 在平面直角坐标系中，点 $P(-3, 2)$ 关于直线 $y=x$ 对称点的坐标是（\quad）

  \optionlist
  \item (-3, -2)
  \item (3, 2)
  \item (2, -3)
  \item (3, -2)
  \end{enumerate}

\item 将 $\triangle ABC$ 的三个顶点的横坐标不变，纵坐标乘以 $-1$，则所得图形（\quad）

  \optionlist
  \item 与原图形关于 $x$ 轴对称
  \item 与原图形关于 $y$ 轴对称
  \item 与原图形关于原点对称
  \item 向 $y$ 轴的负方向平移了一个单位
  \end{enumerate}

\item 平面内点 $A(-1, 2)$ 和点 $B(-1, 6)$ 的对称轴是（\quad）

  \optionlist
  \item $x$ 轴
  \item $y$ 轴
  \item 直线 $y=4$
  \item 直线 $x=-1$
  \end{enumerate}



\item 如图，在 $3 \times 3$ 的正方形网格中由四个格点 $A, B, C, D$，以其中一点为原点，网格线所在直线为坐标轴，建立平面直角坐标系，使其余三个点中存在两个点关于一条坐标轴对称，则原点是（\quad）
  \par\smallskip\noindent
  \begin{minipage}[t]{0.44\linewidth}
  \vspace{0pt}
 \optionlist
  \item $A$ 点
  \item $B$ 点
  \item $C$ 点
  \item $D$ 点
  \end{enumerate}
  \end{minipage}
  \hfill
  \begin{minipage}[t]{0.2\linewidth}
  \questionimage{12.8-5.png}
  \end{minipage}


\item 如图，若 $\triangle A'B'C'$ 与 $\triangle ABC$ 关于直线 $AB$ 对称，则点 $C$ 的对称点 $C'$ 的坐标是（\quad）
  \par\smallskip\noindent
  \begin{minipage}[t]{0.44\linewidth}
  \vspace{0pt}
 \optionlist
  \item (0, -1)
  \item (0, -3)
  \item (3, 0)
  \item (2, 1)
  \end{enumerate}
  \end{minipage}
  \hfill
  \begin{minipage}[t]{0.2\linewidth}
  \questionimage{12.8-6.png}
  \end{minipage}


\item 如图，$\triangle ABC$ 在平面直角坐标系中第二象限内，顶点 $A$ 的坐标是 $(-2, 3)$，先把 $\triangle ABC$ 向右平移 4 个单位得到 $\triangle A_1B_1C_1$，再作 $\triangle A_1B_1C_1$ 关于 $x$ 轴对称图形 $\triangle A_2B_2C_2$，则顶点 $A_2$ 的坐标是（\quad）
  \par\smallskip\noindent
  \begin{minipage}[t]{0.44\linewidth}
  \vspace{0pt}
 \optionlist
  \item (-3, 2)
  \item (2, -3)
  \item (1, -2)
  \item (3, -1)
  \end{enumerate}
  \end{minipage}
  \hfill
  \begin{minipage}[t]{0.2\linewidth}
  \questionimage{12.8-7.png}
  \end{minipage}


\item 如图，若直线 $m$ 经过第二、四象限，且平分坐标轴的夹角，$Rt\triangle AOB$ 与 $Rt\triangle A'OB'$ 关于直线 $m$ 对称，已知 $A(1, 2)$，则点 $A'$ 的坐标为（\quad）
  \par\smallskip\noindent
  \begin{minipage}[t]{0.44\linewidth}
  \vspace{0pt}
 \optionlist
  \item (-1, 2)
  \item (1, -2)
  \item (-1, -2)
  \item (-2, -1)
  \end{enumerate}
  \end{minipage}
  \hfill
  \begin{minipage}[t]{0.25\linewidth}
  \questionimage{12.8-8.png}
  \end{minipage}

\end{enumerate}


\section*{二、填空题（共6小题）}

\blanklist[9]

\item 点 $P(3a+6, 3-a)$ 关于 $x$ 轴的对称点在第四象限内，则 $a$ 的取值范围为\_\_\_\_\_\_。

\item 若点 $P(a, b), Q(c, d)$ 两点关于直线 $x=2$ 对称，则 $a, c$ 间的关系是\_\_\_\_\_\_，$b, d$ 间关系是\_\_\_\_\_\_；若点 $P(a, b), Q(c, d)$ 两点关于直线 $y=-2$ 对称，则 $a, c$ 间的关系是\_\_\_\_\_\_，$b, d$ 间的关系是\_\_\_\_\_\_。

\item 已知 $P_1(a-1, 5)$ 和 $P_2(2, b-1)$ 关于 $x$ 轴对称，则 $(a+b)^{2019}$ 的值为\_\_\_\_\_\_。

\rightquestionimage[width=0.35\textwidth]{0.35\textwidth}{12.8-13.png}

\item 在平面直角坐标系中，点 $A$ 的坐标为 $(-2, 3)$，点 $B$ 的坐标为 $(-1, 6)$。若点 $C$ 与点 $A$ 关于轴对称，则点 $B$ 与点 $C$ 之间的距离为\_\_\_\_\_\_。

\item 在平面直角坐标系 $A$ 中，已知直线 $l: y=x$，作 $A_1(1, 0)$ 关于 $y=x$ 的对称点 $B_1$，将点 $B_1$ 向右水平平移 2 个单位得到点 $A_2$；再作 $A_2$ 关于 $y=x$ 的对称点 $B_2$，将点 $B_2$ 向右水平平移 2 个单位得到点 $A_3$；……按此规律，则点 $B_{2019}$ 的坐标是\_\_\_\_\_\_。

\item 平面直角坐标系中有一点 $A(1, 1)$，对点 $A$ 进行如下操作：

第一步，作点 $A$ 关于 $x$ 轴的对称点 $A_1$，延长线段 $A_1A$ 到点 $A_2$，使得 $2A_1A_2 = A_1A$；

第二步，作点 $A_2$ 关于 $y$ 轴的对称点 $A_3$，延长线段 $A_2A_3$ 到点 $A_4$，使得 $2A_3A_4 = A_2A_3$；

第三步，作点 $A_4$ 关于 $x$ 轴的对称点 $A_5$，延长线段 $A_4A_5$ 到点 $A_6$，使得 $2A_5A_6 = A_4A_5$；

……

则点 $A_2$ 的坐标为\_\_\_\_\_\_，点 $A_{2025}$ 的坐标为\_\_\_\_\_\_。

若点 $A_n$ 的坐标恰好为 $(4^s, 4^t)$（$s, t$ 均为正整数），请写出 $s$ 和 $t$ 的关系式\_\_\_\_\_\_。

\end{enumerate}

\section*{三、解答题（共4小题）}

\solutionlist[15]

\item 已知点 $A(2m+n, 2)$，$B(1, n-m)$，当 $m, n$ 分别为何值时，

(1) $A, B$ 关于 $x$ 轴对称；

(2) $A, B$ 关于 $y$ 轴对称。

\vspace{2cm}

\item 如图，在平面直角坐标系中，$A(-2, 2)$，$B(-3, -2)$。



(1) 若点 $D$ 与点 $A$ 关于 $y$ 轴对称，则点 $D$ 的坐标为\_\_\_\_\_\_。

(2) 将点 $B$ 先向右平移 5 个单位再向上平移 1 个单位得到点 $C$，则点 $C$ 的坐标为\_\_\_\_\_\_。

(3) 求 $A, B, C, D$ 组成的四边形 $ABCD$ 的面积。

\rightquestionimage[width=0.35\textwidth]{0.35\textwidth}{12.8-16.png}

\vspace{3cm}

\item 三角形 $ABC$ 为等腰直角三角形，其中 $\angle A = 90^\circ$，$BC$ 长为 6。

(1) 建立适当的直角坐标系，并写出各个顶点的坐标；

(2) 将 (1) 中各顶点的横坐标都加 2，纵坐标保持不变，与原图案相比，所得的图案有什么变化？

(3) 将 (1) 中各顶点的横坐标不变，将纵坐标都乘 $-1$，与原图案相比，所得的图案有什么变化？

(4) 将 (1) 中各顶点的横坐标都乘 $-2$，纵坐标保持不变，与原图案相比，所得的图案有什么变化？

\vspace{7cm}

\item 如图所示，$\triangle COB$ 是由 $\triangle AOB$ 经过某种变换后得到的图形，观察点 $A$ 与点 $C$ 的坐标之间的关系，解答下列问题：

% 插入图片：此处应有第18题图


(1) 若点 $M$ 的坐标为 $(x, y)$，则它的对应点 $N$ 的坐标为\_\_\_\_\_\_。

(2) 若点 $P(a, 2)$ 与点 $Q(-3, b)$ 关于 $x$ 轴对称，求代数式
\[\frac{1}{ab} + \frac{1}{(a-1)(b-1)} + \frac{1}{(a-2)(b-2)} + \cdots + \frac{1}{(a-10)(b-10)}\]
的值。

\rightquestionimage[width=0.35\textwidth]{0.35\textwidth}{12.8-18.png}


\end{enumerate}


\chapter{12.9 等腰三角形}



\section*{三、解答题（共4小题）}

\solutionlist[14]


\item 阅读下列材料，按要求解答问题：

如图 1，在 $\triangle ABC$ 中，$\angle A = 2\angle B$，且 $\angle A = 60^\circ$。小明通过以下计算：由题意，$\angle B = 30^\circ$，$\angle C = 90^\circ$，$c = 2b$，$a = \sqrt{3}b$，得 $a^2 - b^2 = (\sqrt{3}b)^2 - b^2 = 2b^2 = b \cdot c$。即 $a^2 - b^2 = bc$。于是，小明猜测：对于任意的 $\triangle ABC$，当 $\angle A = 2\angle B$ 时，关系式 $a^2 - b^2 = bc$ 都成立。

(1) 如图 2，请你用以上小明的办法，对等腰直角三角形进行验证，判断小明的猜测是否正确，并写出验证过程。

\rightquestionimage[width=0.5\textwidth]{0.5\textwidth}{12.9-14-1.png}

\vspace{3cm}

(2) 如图 3，$\angle A = 30^\circ$，$\angle B = 15^\circ$，试设计方法，判断小明的猜测是否正确，并写出验证过程。

\rightquestionimage[width=0.5\textwidth]{0.5\textwidth}{12.9-14-2.png}

\vspace{3cm}

(3) 若一个三角形的三边长恰为三个连续偶数，且 $\angle A = 2\angle B$，请直接写出这个三角形三边的长不必说明理由。

\vspace{3cm}



\item 如图，在平面直角坐标系中，已知点 $A(0, 8)$，$B(6, 0)$。以 $\triangle AOB$ 的一边为边画等腰三角形，使它的第三个顶点在 $\triangle AOB$ 的一边上。请在图①、图②中分别画出一个符合条件的等腰三角形，且两个图中的等腰三角形不相同，在图中标明所画等腰三角形的第三个顶点的坐标（不要求尺规作图，不写求解过程）。

% 插入图片：此处应有图①

% 插入图片：此处应有图②

\vspace{6cm}



\item 如图，在 $\triangle ABC$ 中，$AB = AC = a$，$BC = b$，且 $2a > b$，$BG \perp AC$ 于 $G$，$DE \perp AB$ 于 $E$，$DF \perp AC$ 于 $F$。

(1) 在图 (1) 中，$D$ 是 $BC$ 边上的中点，计算 $DE + DF$ 和 $BG$ 的长（用 $a, b$ 表示），并判断 $DE + DF$ 与 $BG$ 的关系。


(2) 在图 (2) 中，$D$ 是线段 $BC$ 上的任意一点，$DE + DF$ 与 $BG$ 的关系是否仍然成立？如果成立，证明你的结论；如果不成立，请说明理由。


(3) 在图 (3) 中，$D$ 是线段 $BC$ 延长线上的点，探究 $DE, DF$ 与 $BG$ 的关系。（不要求证明）

\centerquestionimage[width=0.7\textwidth]{0.7\textwidth}{12.9-16.png}

\vspace{5cm}


\newpage

\item 如图所示，在 $\triangle ABC$ 中，$AB = AC$，$\angle BAC = 80^\circ$，$P$ 在 $\triangle ABC$ 内，$\angle PBC = 10^\circ$，$\angle PCB = 30^\circ$，求 $\angle PAB$ 的度数。

\rightquestionimage[width=0.5\textwidth]{0.5\textwidth}{12.9-17.png}

\vspace{3cm}

\end{enumerate}


\chapter{12.10-2 等腰三角形判定}

\section*{一、选择题（共 8 小题）}

\questionlist[8]
\item
  \begin{questionwithimage}{12.10-1-8.png}
  \item 在 $4 \times 4$ 的方格表中给出如图所示的 8 个点，任选三个作为三角形的顶点，共可构成（\quad）个等腰三角形。
  \optionlist
  \item 5
  \item 6
  \item 7
  \item 8
  \end{enumerate}
  \end{questionwithimage}


\end{enumerate}

\section*{二、填空题（共 6 小题）}

\blanklist[9]

\item 已知 $a, b, c$ 为三个正整数，如果 $a + b + c = 12$，那么以 $a, b, c$ 为边能组成的三角形是：①等腰三角形；②等边三角形；③直角三角形；④钝角三角形。以上符合条件的正确结论是\_\_\_\_\_\_。（只填序号）

\item 如图，$AD$ 是 $\triangle ABC$ 的边 $BC$ 上的高，现给出下列条件：① $\angle BAD = \angle ACD$；② $\angle BAD = \angle CAD$；③ $BD = CD$；④ $AB + BD = AC + CD$，若添加这些条件中的某一个就能推出 $\triangle ABC$ 是等腰三角形，这个条件可以是\_\_\_\_\_\_。（只填序号）

\rightquestionimage[width=0.35\textwidth]{0.35\textwidth}{12.10-11.png}

\item 如图，$AD$ 是直角三角形 $\triangle ABC$ 斜边上的中线，把 $ADC$ 沿 $AD$ 对折，点 $C$ 落在点 $C'$ 处，连接 $CC'$，则图中共有等腰三角形\_\_\_\_\_\_个。

\rightquestionimage[width=0.35\textwidth]{0.35\textwidth}{12.10-12.png}

\item 如图，组成虚线网格的每个小正三角形的边长都为 1，若有格点（小正三角形的顶点）$C$，使 $\triangle ABC$ 为等腰三角形，那么这样的 $C$ 点共有\_\_\_\_\_\_个。

\rightquestionimage[width=0.35\textwidth]{0.35\textwidth}{12.10-13.png}

\item 以原点 $O(0, 0)$，$A(1, 2)$ 为顶点的 $\triangle AOB$ 是等腰三角形，且点 $B$ 在坐标轴上，满足条件的点 $B$ 有\_\_\_\_\_\_个，其中横纵坐标均为整数的坐标为\_\_\_\_\_\_。

\item 如图，点 $A$ 的坐标为 $(0, 8)$，点 $B$ 的坐标为 $(6, 0)$，在 $x$ 轴上确定一点 $P$，使 $\triangle PAB$ 为一个等腰三角形，则 $R$ 点的坐标可以是\_\_\_\_\_\_。

\rightquestionimage[width=0.35\textwidth]{0.35\textwidth}{12.10-14.png}

\end{enumerate}

\section*{三、解答题（共 4 小题）}

\solutionlist[15]

\item 如图：已知 $AB = AE$，$BC = ED$，$\angle B = \angle E$，$AF \perp CD$，$F$ 为垂足，求证：① $AC = AD$；② $CF = DF$。

\rightquestionimage[width=0.35\textwidth]{0.35\textwidth}{12.10-15.png}

\vspace{2cm}

\item 已知，如图，在 $\triangle ABC$ 中，$AD \perp BC$，垂足为点 $D$，$BE \perp AC$，垂足为点 $E$，$M$ 为 $AB$ 边的中点，连接 $ME$、$MD$、$ED$。

(1) 求证：$\triangle MED$ 为等腰三角形；

(2) 求证：$\angle EMD = 2\angle DAC$。

\rightquestionimage[width=0.35\textwidth]{0.35\textwidth}{12.10-16.png}

\vspace{5cm}

\item 已知：一张直角三角形纸片如图1放置在平面直角坐标系中，一条直角边 $OA$ 落在 $x$ 轴正半轴上，另一条直角边 $OB$ 落在 $y$ 轴正半轴上，且 $OA = 8$，$OB = 6$。现再找一个与 $Rt\triangle ABO$ 有一条公共边且不重叠的三角形，使它们拼在一起后能构成一个大的等腰三角形。例如：如图2，$\triangle CBO$ 与 $\triangle ABO$ 拼成等腰 $\triangle ABC$，则点 $C$ 坐标为 $(-2, 0)$。请直接写出除图2情况外，其他所有的所拼成的等腰三角形中除 $A$、$B$、$O$ 三点外另一顶点 $P$ 的坐标。

\rightquestionimage[width=0.50\textwidth]{0.50\textwidth}{12.10-17.png}

\vspace{5cm}

\end{enumerate}

\chapter{12.10-2 等腰三角形的性质与判定}


\section*{二、填空题（共6小题）}

\blanklist[9]

\item 如图，在 $\triangle ABC$ 中，$\angle BAC = 150^\circ$，$AD \perp BC$ 于 $D$，且 $AB + BD = DC$，那么 $\angle C =$ \_\_\_\_\_\_。

\rightquestionimage[width=0.35\textwidth]{0.35\textwidth}{12.10-2-9.png}

\item 如图，$\triangle ABC$ 中，$AB = AC$，$BD$ 平分 $\angle ABC$ 交 $AC$ 于 $G$，$DM \parallel BC$ 交 $\angle ABC$ 的外角平分线于 $M$，交 $AB$、$AC$ 于 $F$、$E$，下列结论：

① $MB \perp BD$；② $FD = EC$；③ $EC = EF + DG$；④ $CE = \frac{1}{2}MD$。

其中正确的有\_\_\_\_\_\_（填序号）。

\rightquestionimage[width=0.35\textwidth]{0.35\textwidth}{12.10-2-10.png}

\item 如图：已知 $\angle BAD = \angle DAC = 90^\circ$，$AD \perp AE$，且 $AB + AC = BE$，则 $\angle B =$ \_\_\_\_\_\_。

\rightquestionimage[width=0.35\textwidth]{0.35\textwidth}{12.10-2-11.png}

\item 等腰三角形一腰上的高等于该三角形某一条边的长度的一半，则其顶角可能为\_\_\_\_\_\_。

\item 如图，点 $B$ 是线段 $AC$ 的中点，过点 $C$ 的射线 $CE$ 与 $AC$ 成 $60^\circ$ 的角，点 $P$ 为射线 $CE$ 上一动点，给出以下四个结论：

① 当 $AP \perp CE$，垂足为 $P$ 时，$\angle APB = 30^\circ$；

② 当 $CP = AC$ 时，$\angle APB = 30^\circ$；

③ 在射线 $CE$ 上，使 $\triangle APC$ 为直角三角形的点 $P$ 只有 1 个；

④ 在射线 $CE$ 上，使 $\triangle APC$ 为等腰三角形的点 $P$ 只有 1 个；

其中正确结论的序号是\_\_\_\_\_\_。

\rightquestionimage[width=0.35\textwidth]{0.35\textwidth}{12.10-2-13.png}

\item 如图，$\triangle ABC$ 中，$AB = BC = AC = 12\text{cm}$，现有两点 $M$、$N$ 分别从点 $A$、$B$ 同时出发，沿三角形的边运动，已知点 $M$ 的速度为 $1\text{cm/s}$，点 $N$ 的速度为 $2\text{cm/s}$。当点 $N$ 第一次到达 $B$ 点时，$M$、$N$ 同时停止运动。


(1) 当点 $M$、$N$ 运动\_\_\_\_\_\_秒时，$M$、$N$ 两点重合；

(2) 当点 $M$、$N$ 运动\_\_\_\_\_\_秒后，$M$、$N$ 与 $\triangle ABC$ 中的某个顶点可得到等腰三角形。

\rightquestionimage[width=0.35\textwidth]{0.35\textwidth}{12.10-2-14.png}

\end{enumerate}

\section*{三、解答题（共3小题）}

\solutionlist[15]

\item 已知：如图，$AD$ 是 $\triangle BAC$ 的角平分线，$BE \perp AD$ 交 $AD$ 的延长线于点 $E$，$EF \parallel AC$ 交 $AB$ 于点 $F$，$AE = 8$，$EF = 5$。求 $BE$ 的长。

\rightquestionimage[width=0.35\textwidth]{0.35\textwidth}{12.10-2-15.png}

\newpage

\item 在 $\triangle ABC$ 中，$\angle C = 90^\circ$，$AC = BC = 2$，将一块三角板的直角顶点放在斜边 $AB$ 的中点 $P$ 处，将此三角板绕点 $P$ 旋转，三角板的两直角边分别交射线 $AC$、$CB$ 与点 $D$、点 $E$，图①、②、③是旋转得到的三种图形。

(1) 观察线段 $PD$ 和 $PE$ 之间的有怎样的关系，并以图②为例，加以说明；

(2) $\triangle PBE$ 是否构成等腰三角形？若能，指出所有的情况（即求出 $\triangle PBE$ 为等腰三角形时 $CE$ 的长）；若不能请说明理由。

\centerquestionimage[width=0.70\textwidth]{0.70\textwidth}{12.10-2-16.png}

\vspace{5cm}

\item 如图，在 $\triangle ABC$ 中，$AB = AC$，$P$ 为 $\triangle ABC$ 内一点，且 $\angle BAP = 70^\circ$，$\angle ABP = 40^\circ$。


(1) 求证：$\triangle ABP$ 是等腰三角形；

(2) 连接 $PC$，当 $\angle PCB = 30^\circ$ 时，求 $\angle PBC$ 的度数。

\rightquestionimage[width=0.40\textwidth]{0.40\textwidth}{12.10-2-17.png}

\end{enumerate}


\chapter{12.11-1 等腰三角形性质}


\section*{二、填空题（共6小题）}


\blanklist[7]

\item 如图，阴影部分是边长为1的小正三角形，$A, B, C, D, E, F, G, H$ 分别是8个正三角形，则 $A$ 和 $B$ 的边长分别是（\quad ）
  \par\smallskip\noindent
  \begin{minipage}[t]{0.44\linewidth}
  \vspace{0pt}
 \optionlist
  \item 2, 4
  \item 2.5, 5
  \item 3, 6
  \item 4, 8
  \end{enumerate}
  \end{minipage}
  \hfill
  \begin{minipage}[t]{0.2\linewidth}
  \questionimage{12.11-7.png}
  \end{minipage}


\item 图①是一块边长为1，周长记为 $P_1$ 的正三角形纸板，沿图①的底边剪去一块边长为 $\frac{1}{2}$ 的正三角形纸板后得到图②，然后沿同一底边依次剪去一块更小的正三角形纸板（即其边长为前一块被剪掉正三角形纸板边长的 $\frac{1}{2}$）后，得图③，④，…，记第 $n$（$n \geq 3$）块纸板的周长为 $P_n$，则 $P_n - P_{n-1}$ 的值为（    ）
  \par\smallskip\noindent
  \begin{minipage}[t]{0.44\linewidth}
  \vspace{0pt}
  \optionlist
  \item $\left( \frac{1}{4} \right)^{n-1}$
  \item $\left( \frac{1}{4} \right)^n$
  \item $\left( \frac{1}{2} \right)^{n-1}$
  \item $\left( \frac{1}{2} \right)^n$
  \end{enumerate}
    \end{minipage}
  \hfill
  \begin{minipage}[t]{0.55\linewidth}
  \questionimage{12.11-8.png}
  \end{minipage}

% 插入图片：此处应有图①、②、③、④

\end{enumerate}

\section*{二、填空题（共6小题）}

\blanklist[9]

\item 如图，已知 $\triangle ABC$ 是等边三角形，点 $O$ 是 $BC$ 上任意一点，$OE, OF$ 分别与两边垂直，等边三角形的高为5，则 $OE + OF$ 的值为\_\_\_\_\_\_。

\rightquestionimage[width=0.35\textwidth]{0.35\textwidth}{12.11-9.png}

\item 三个等边三角形的位置如图所示，若 $\angle 3 = 40^\circ$，则 $\angle 1 + \angle 2 =$\_\_\_\_\_\_。

\rightquestionimage[width=0.35\textwidth]{0.35\textwidth}{12.11-10.png}

\item 如图，$\triangle ABC$ 是边长为3的等边三角形，$\triangle BDC$ 是等腰三角形，且 $\angle BDC = 120^\circ$。以 $D$ 为顶点作一个 $60^\circ$ 角，使其两边分别交 $AB$ 于点 $M$，交 $AC$ 于点 $N$，连接 $MN$，则 $\triangle AMN$ 的周长为\_\_\_\_\_\_。

\rightquestionimage[width=0.35\textwidth]{0.35\textwidth}{12.11-11.png}

\item 将一个平面图形分成面积相等的两部分的直线叫做该平面图形的“面线”，“面线”被这个平面图形截得的线段叫做该图形的“面径”，例如圆的直径就是它的“面径”。已知等边三角形的边长为2，则它的“面径”长可以是\_\_\_\_\_\_（写出2个）。


\item 如图，将边长为1的正三角形 $OAP$ 沿 $x$ 轴正方向连续无滑动的翻转2019次，点 $P$ 依次落在点 $P_1, P_2, \ldots, P_{2019}$ 的位置，则点 $P_{2019}$ 的横坐标为\_\_\_\_\_\_。

\rightquestionimage[width=0.35\textwidth]{0.35\textwidth}{12.11-13.png}

\item 如图，一只蜘蛛结了一张很大的网，直线 $MN, PQ$ 与 $x$ 轴所夹的锐角都为 $60^\circ$，第1个结点 $A_1$ 在原点，此后各个结点均按逆时针排列，同一直线上相邻两个结点之间的距离都是1个单位长，那么第91个结点 $A_{91}$ 的坐标为\_\_\_\_\_\_。

\rightquestionimage[width=0.35\textwidth]{0.35\textwidth}{12.11-14.png}

\end{enumerate}

\newpage

\section*{三、解答题（共3小题）}

\solutionlist[15]

\item 如图，以 $\triangle ABC$ 的两边 $AB, AC$ 向外作等边三角形 $ABE$ 和等边三角形 $ACD$，连接 $BD, CE$，交于 $O$。

(1) 试写出图中和 $BD$ 相等的一条线段并说明你的理由；

(2) 求出 $BD$ 和 $CE$ 的夹角大小，若改变 $\triangle ABC$ 的形状，这个夹角的度数会发生变化吗？请说明理由。

\rightquestionimage[width=0.35\textwidth]{0.35\textwidth}{12.11-15.png}

\vspace{3cm}

\item 将两个等边 $\triangle ABC$ 和 $\triangle DEF$（$DE > AB$）如图所示摆放，点 $D$ 是 $BC$ 上的一点（除 $B, C$ 点外）。把 $\triangle DEF$ 绕顶点 $D$ 顺时针旋转一定的角度，使得边 $DE, DF$ 与 $\triangle ABC$ 的边（除 $BC$ 边外）分别相交于点 $M, N$。

(1) $\angle BMD$ 和 $\angle CDN$ 相等吗？

(2) 画出使 $\angle BMD$ 和 $\angle CDN$ 相等的所有情况的图形；

(3) 在 (2) 题中任选一种图形说明 $\angle BMD$ 和 $\angle CDN$ 相等的理由。


\rightquestionimage[width=0.35\textwidth]{0.35\textwidth}{12.11-16.png}

\end{enumerate}

\newpage

\solutionlist[17]

\item 已知：等边三角形 $ABC$，

(1) 如图1，$P$ 为等边 $\triangle ABC$ 外一点，且 $\angle BPC = 120^\circ$。试猜想线段 $BP$、$PC$、$AP$ 之间的数量关系，并证明你的猜想；

(2) 如图2，$P$ 为等边 $\triangle ABC$ 内一点，且 $\angle APD = 120^\circ$。求证：$PA + PD + PC > BD$。

\centerquestionimage[width=0.70\textwidth]{0.70\textwidth}{12.11-17.png}


\end{enumerate}

\chapter{12.11-2 等边三角形的判定}



\section*{12.11.2 等边三角形的判定}

\subsection*{一、选择题（共 8 小题）}

\questionlist

\item 已知 $a, b, c$ 是 $\triangle ABC$ 的三边，且 $a^2 + b^2 + c^2 = ab + ac + bc$，则 $\triangle ABC$ 是（\quad）

  \optionlist
  \item 等腰三角形
  \item 直角三角形
  \item 等边三角形
  \item 等腰直角三角形
  \end{enumerate}

\item 下列三角形：

\begin{enumerate}[label=(\arabic*)]
\item 有两个角等于 $60^\circ$；
\item 有一个角等于 $60^\circ$ 的等腰三角形；
\item 三个外角（每个顶点处各取一个外角）都相等的三角形；
\item 一腰上的中线也是这条腰上的高的等腰三角形。
\end{enumerate}

其中是等边三角形的有（\quad）

\optionlist
\item 1,2,3
\item 1,2,4
\item 1,3
\item 1,2,3,4
\end{enumerate}

\item
  \begin{questionwithimage}{12-11-2-3.jpg}
  如图，$E$ 是等边 $\triangle ABC$ 中 $AC$ 边上的点，$\angle 1 = \angle 2$，$BE = CD$，则 $\triangle ADE$ 的形状是（\quad）
  \optionlist
  \item 等腰三角形
  \item 等边三角形
  \item 不等边三角形
  \item 不能确定形状
  \end{enumerate}
  \end{questionwithimage}

\item 下列推理错误的是（\quad）

\optionlist
\item 在 $\triangle ABC$ 中，$\because \angle A = \angle B = \angle C$，$\therefore \triangle ABC$ 为等边三角形
\item 在 $\triangle ABC$ 中，$AB = AC$，且 $\angle B = \angle C$，$\triangle ABC$ 为等边三角形
\item 在 $\triangle ABC$ 中，$\angle A = 60^\circ$，$\angle B = 60^\circ$，$\triangle ABC$ 为等边三角形
\item 在 $\triangle ABC$ 中，$AB = AC$，$\angle B = 60^\circ$，$\triangle ABC$ 为等边三角形
\end{enumerate}

\item 三角形中有两条中线分别平分它的两个内角，则这个三角形是（\quad）

  \optionlist
  \item 直角三角形
  \item 等腰三角形
  \item 等边三角形
  \item 等腰直角三角形
  \end{enumerate}

\item
  \begin{questionwithimage}{12-11-2-6.png}
  如图，等边 $\triangle DEF$ 的顶点分别在等边 $\triangle ABC$ 的各边上，且 $DE \perp BC$ 于 $E$，若 $AB = 1$，则 $DB$ 的长为（\quad）
  \optionlist
  \item $\frac{1}{2}$
  \item $\frac{1}{3}$
  \item $\frac{2}{3}$
  \item $\frac{3}{4}$
  \end{enumerate}
  \end{questionwithimage}

\item
  \begin{questionwithimage}{12-11-2-7.png}
  如图，在等边三角形 $ABC$ 中，$BC = 2$，$D$ 是 $AB$ 的中点，过点 $D$ 作 $DF \perp AC$ 于点 $F$，过点 $F$ 作 $EF \perp BC$ 于点 $E$，则 $BE$ 的长为（\quad）
  \optionlist
  \item 1
  \item $\frac{3}{2}$
  \item $\frac{5}{4}$
  \item $\frac{4}{3}$
  \end{enumerate}
  \end{questionwithimage}

\item 如图，$\angle MON = 30^\circ$。点 $A_1, A_2, A_3, \ldots$ 在射线 $ON$ 上，点 $B_1, B_2, B_3, \ldots$ 在射线 $OM$ 上。$\triangle A_1B_1A_2$、$\triangle A_2B_2A_3$、$\triangle A_3B_3A_4$、$\ldots$ 均为等边三角形。若 $OA_1 = 1$，则 $\triangle A_7B_7A_8$ 的边长为（\quad）
  \begin{questionwithimage}{12-11-2-8.png}
  \optionlist
  \item 6
  \item 12
  \item 32
  \item 64
  \end{enumerate}
  \end{questionwithimage}

\end{enumerate}

\section*{二、填空题（共 6 小题）}

\blanklist[9]
\item
  若正三角形的边长为 $2\sqrt{5}\,cm$，则这个正三角形的高是\underline{\hspace{2cm}}

\item
  在平面直角坐标系中，等边$\triangle ABC$的顶点$A$（-6，0)，$B$（2，0)，则顶点$C$的坐标为\underline{\hspace{2cm}}

\item
  $\triangle ABC$ 是等边三角形，点$D$是$BC$边上任意一点， $DE \perp AB$ 于点$E$， $DF \perp AC$ 于点$F$. 若 $BC=$ 4，则 $BE+CF=$\underline{\hspace{2cm}}

\item
  已知a、b、c是 $\triangle ABC$ 的三边的长，且满足 $a^{2}+2b^{2}+c^{2}-2b(a+c)=0$，则此三角形的形状为\underline{\hspace{2cm}}

\item
  如图， $\triangle ABC$ 是一个边长为1的等边三角形， $BB_{1}$ 是 $\triangle ABC$ 的高，$B_{1}B_{2}$ 是 $\triangle ABB_{1}$ 的高 $B_{2}B_{3}$ 是 $\triangle AB_{1}B_{2}$ 的高， $\ldots , \quad B_{n-1}B_{n}$ 是 $\triangle AB_{n-2}B_{n-1}$ 的高，$B_{n-1}B_{n}$ 的长是\underline{\hspace{2cm}}

  \centerquestionimage{0.4\textwidth}{12-11-3-1.jpg}

\item
  在 $\triangle ABC$ 中， $\angle C=90^{\circ}$ $AB=2\,cm$ $BC=1\,cm$ .以点$C$为顶点作一个等边三角形，使其他两个顶点在 $\triangle ABC$ 的边上，则这个等边三角形的面积为\underline{\hspace{2cm}}

\end{enumerate}

\section*{三、解答题（共 1 小题）}

\solutionlist[15]
\item
  在 $\mathrm{Rt}\triangle ABC$ 中， $\angle ACB=90^{\circ}$，BD是 $\triangle ABC$ 的角平分线.
(1) 如图1，若 $\scriptstyle AD=BD$，求 $\angle A$ 的度数；
(2)如图2，在（1)的条件下，作 $DE \perp AB$ 于E，连接 $EC$. 求证： $\triangle EBC$ 是等边三角形.

  \begin{questionimagegroup}
  \questionimageitem{0.32\textwidth}{12-11-3-2.jpg}{图1}
  \questionimagegap
  \questionimageitem{0.32\textwidth}{12-11-3-3.jpg}{图2}
  \end{questionimagegroup}

\item
  如图，在 $\triangle ABC$ 中，$\angle A=60^{\circ}$，$BE \perp AC$，垂足为 $E$，$CF \perp AB$，垂足为 $F$，点 $D$ 是 $BC$ 的中点，$BE$，$CF$ 交于点 $M$。

  (1) 如果 $AB=AC$，求证：$\triangle DEF$ 是等边三角形；

  (2) 如果 $AB\ne AC$，试猜想 $\triangle DEF$ 是不是等边三角形？如果 $\triangle DEF$ 是等边三角形，请加以证明；如果 $\triangle DEF$ 不是等边三角形，请说明理由；

  (3) 如果 $CM=4$，$FM=5$，求 $BE$ 的长度。

  \begin{questionimagegroup}
    \questionimageitem[width=\linewidth,height=0.22\textheight,keepaspectratio]{0.32\textwidth}{12-11-2-16-1.png}{}
  \questionimagegap[0.04\textwidth]
    \questionimageitem[width=\linewidth,height=0.22\textheight,keepaspectratio]{0.32\textwidth}{12-11-2-16-2.png}{}
  \end{questionimagegroup}

\vspace{5cm}

\item
  如图所示，$\triangle OAB$ 是边长为 $2+\sqrt{3}$ 的等边三角形，其中 $O$ 是坐标原点，顶点 $A$ 在 $x$ 轴的正方向上。将 $\triangle OAB$ 折叠，使点 $B$ 落在边 $OA$ 上，记为 $B'$，折痕为 $EF$。

  (1) 设 $OB'$ 的长为 $x$，$\triangle OB'E$ 的周长为 $c$，用含有 $x$ 的式子表示 $c$；

  (2) 当 $B'E \parallel y$ 轴时，求点 $B'$ 和点 $E$ 的坐标；

  (3) 当 $B'$ 在 $OA$ 上运动但不与 $O$、$A$ 重合时，能否使 $\triangle EB'F$ 成为直角三角形？若能，请求出点 $B'$ 的坐标；若不能，请说明理由。

\centerquestionimage[width=0.35\textwidth]{0.35\textwidth}{12-11-2-17-1.png}

\vspace{5cm}

\item 已知 $\triangle ABC$ 是等腰三角形，$AB = AC$，$D$ 为边 $BC$ 上任意一点，$DE \perp AB$ 于 $E$，$DF \perp AC$ 于 $F$，且 $E, F$ 分别在边 $AB, AC$ 上。

(1) 如图 a，当 $\triangle ABC$ 是等边三角形时，证明：$AE + AF = \frac{3}{2}BC$。

(2) 如图 b，若 $\triangle ABC$ 中，$\angle BAC = 120^\circ$，探究线段 $AE, AF, AB$ 之间的数量关系，并对你的猜想加以证明。

(3) 如图 c，若 $\triangle ABC$ 中，$AB = 10, BC = 16, EF = 6$，利用你对 (1), (2) 两题的解题思路计算出线段 $CD$（$BD > CD$）的长。

\centerquestionimage[width=0.60\textwidth]{0.60\textwidth}{12-11-2-18.png}


\end{enumerate}


\chapter{12.11-3等边三角形综合}

\section*{一、选择题（共 9 小题）}

\questionlist[1]
\item
  关于等边三角形，下列说法中错误的是（\quad）
  \optionlist
  \item 等边三角形中，各边都相等
  \item 等边三角形是特殊的等腰三角形
  \item 三个角都等于 $60^{\circ}$ 的三角形是等边三角形
  \item 有一个角为 $60^{\circ}$ 的等腰三角形不是等边三角形
  \end{enumerate}

\item
  如图, $\triangle MNP$ 中, $\angle P=60^{\circ}$ , $MN=NP$ , $MQ \perp PN$ , 垂足为 $Q$ , 延长 $MN$ 至 $G$ , 取 $NG=NQ$ , 若 $\triangle MNP$ 的周长为 12, $MQ=a$ , 则 $\triangle MGQ$ 周长是 （\quad）
  \begin{questionwithimage}{2026-06-10-14-11_11-1.jpg}
  \optionlist
  \item $8+2a$
  \item $8+a$
  \item $6+a$
  \item $6+2a$
  \end{enumerate}
  \end{questionwithimage}

\item
  如图, 已知 $\angle ABC=120^{\circ}$ , $BD$ 平分 $\angle ABC$ , $\angle DAC=60^{\circ}$ , 若 $AB=2$ , $BC=3$ , 则 $BD$ 的长是 （\quad）
  \begin{questionwithimage}{2026-06-10-14-11_11-2.jpg}
  \optionlist
  \item 5
  \item 7
  \item 8
  \item 9
  \end{enumerate}
  \end{questionwithimage}

\item
  正三角形 $\triangle ABC$ 的边长为 3，依次在边 $AB$、$BC$、$CA$ 上取点 $A_{1}$、$B_{1}$、$C_{1}$，使 $AA_{1}=BB_{1}=CC_{1}=1$，则 $\triangle A_{1}B_{1}C_{1}$ 的面积是（\quad）
  \optionlist
  \item $\frac{\sqrt{3}}{4}$
  \item $\frac{3\sqrt{3}}{4}$
  \item $\frac{9}{4}$
  \item $\frac{9\sqrt{3}}{4}$
  \end{enumerate}

\item
  如图, $\triangle ABC$ 中, $\angle BAC=120^{\circ}$ , $AD \perp BC$ 于 $D$ , 且 $AB+BD=DC$ , 则 $\angle C$ 的大小是 （\quad）
  \begin{questionwithimage}{2026-06-10-14-11_11-3.jpg}
  \optionlist
  \item $20^{\circ}$
  \item $25^{\circ}$
  \item $30^{\circ}$
  \item 大于 $30^{\circ}$
  \end{enumerate}
  \end{questionwithimage}

\item
  如图,在 $Rt\triangle ABC$ , $\angle ACB=90^{\circ}$ , $CD$、 $CE$ 是斜边上的高和中线, $AC=CE=10cm$ , 则 $BD$ 长为 （\quad）
  \begin{questionwithimage}{2026-06-10-14-11_11-4.jpg}
  \optionlist
  \item $5cm$
  \item $10cm$
  \item $15cm$
  \item $25cm$ ,
  \end{enumerate}
  \end{questionwithimage}

\item
  如图, 六边形 $ABCDEF$ 中, 每一个内角都是 $120^{\circ}, AB=12, BC=30, CD=8, DE=28$ . 求这个六边形的周长为 （\quad）
  \begin{questionwithimage}{2026-06-10-14-11_11-5.jpg}
  \optionlist
  \item 125
  \item 126
  \item 116
  \item 108
  \end{enumerate}
  \end{questionwithimage}

\item
  如图, $P$ 是等边三角形 $ABC$ 内的一点, 且 $PA=3$ , $PB=4$ , $PC=5$ , 以 $BC$ 为边在 $\triangle ABC$ 外作 $\triangle BQC\cong \triangle BPA$ , 连接 $PQ$ , 则以下结论错误的是 （\quad）
  \begin{questionwithimage}{2026-06-10-14-11_11-6.jpg}
  \optionlist
  \item $\triangle BPQ$ 是等边三角形
  \item $\triangle PCQ$ 是直角三角形
  \item $\angle APB=150^\circ$
  \item $\angle APC=135^\circ$
  \end{enumerate}
  \end{questionwithimage}

\item
  如图所示, 在 $\triangle ABC$ 中, $AB=AC$ , $D$、 $E$ 是 $\triangle ABC$ 内两点, $AD$ 平分 $\angle BAC$ . $\angle EBC=\angle E=60^{\circ}$ , 若 $BE=6$ , $DE=2$ , 则 $BC$ 的长度是 （\quad）
  \begin{questionwithimage}{mineru-1.jpg}
  \optionlist
  \item 6
  \item 8
  \item 9
  \item 10
  \end{enumerate}
  \end{questionwithimage}

\end{enumerate}

\section*{二、填空题（共 5 小题）}

\blanklist[10]
\item
  如图, 在一个正方体的 2 个面上画了两条对角线 $AB$、 $AC$ , 那么这两条对角线的夹角是 \_\_\_\_.\underline{\hspace{2cm}}

  \centerquestionimage{0.35\textwidth}{mineru-2.jpg}

\item
  由于木质衣架没有柔性，在挂置衣服的时候不太方便操作。小敏设计了一种衣架，在使用时能轻易收拢，然后套进衣服后松开即可。如图\circled{1}，衣架杆 $OA=OB=18cm$，若衣架收拢时， $\angle AOB=60^{\circ}$，如图\circled{2}，则此时 $A, B$ 两点之间的距离是 \_\_\_\_ cm。\underline{\hspace{2cm}}

  \begin{questionimagegroup}
  \questionimageitem{0.32\textwidth}{mineru-3.jpg}{图\circled{1}}
  \questionimagegap
  \questionimageitem{0.32\textwidth}{mineru-4.jpg}{图\circled{2}}
  \end{questionimagegroup}

\item
  如图, $\triangle ABC$ 是等边三角形, 高 $AD$、$BE$ 相交于点 $H$ , $BC=4\sqrt{3}$ , 在 $BE$ 上截取 $BG=2$ , 以 $GE$ 为边作等边三角形 $GEF$ , 则 $\triangle ABH$ 与 $\triangle GEF$ 重叠 (阴影) 部分的面积为 \_\_\_\_.\underline{\hspace{2cm}}

  \centerquestionimage{0.4\textwidth}{mineru-5.jpg}

\item
  如图所示, 已知 $\angle 1=\angle 2$ , $AD=BD=4$ , $CE \perp AD$ , $2CE=AC$ , 那么 $DE$ 的长是 \_\_\_\_.\underline{\hspace{2cm}}

  \centerquestionimage{0.4\textwidth}{mineru-6.jpg}

\item
  如图, 直线 $a \parallel b$ , $\triangle ABC$ 是等边三角形, 点 $A$ 在直线 $a$ 上, 边 $BC$ 在直线 $b$ 上, 把 $\triangle ABC$ 沿 $BC$ 方向平移 $BC$ 的一半得到 $\triangle A'B'C'$ (如图\circled{1}); 继续以上的平移得到图\circled{2}, 再继续以上的平移得到图\circled{3}, ...; 请问在第 100 个图形中等边三角形的个数是 \_\_\_\_.\underline{\hspace{2cm}}

  \begin{questionimagegroup}
  \questionimageitem{0.28\textwidth}{mineru-7.jpg}{\circled{1}}
  \questionimagegap
  \questionimageitem{0.28\textwidth}{mineru-8.jpg}{\circled{2}}
  \questionimagegap
  \questionimageitem{0.28\textwidth}{mineru-9.jpg}{\circled{3}}
  \end{questionimagegroup}

\end{enumerate}

\section*{三、解答题（共 1 小题）}

\solutionlist[15]
\item
  如图, 已知 $\triangle ABC$ 为等边三角形, 延长 $BC$ 到 $D$ , 延长 $BA$ 到 $E$ , 并且使 $AE=BD$ , 连接 $CE$ , $DE$ . 求证: $EC=ED$ .

  \centerquestionimage{0.4\textwidth}{mineru-10.jpg}

\end{enumerate}



\end{document}
